% Original debugging lab. Arithmetic examples assume IEEE binary64, round-to-nearest.
\documentclass[14pt,aspectratio=1610]{beamer}
\usetheme{Rochester}
\usecolortheme{beaver}
\usepackage{amsmath,listings,alltt}
\setbeamertemplate{navigation symbols}{}
\lstset{basicstyle=\ttfamily\footnotesize,columns=fullflexible,keepspaces=true}
\title{A plausible number can still be wrong}
\subtitle{A numerical debugging lab}
\author{Scientific computing studio}
\date{Lab 12}
\begin{document}
\begin{frame}[plain]\titlepage\end{frame}
\begin{frame}[fragile]{Run this before proposing a fix}
\begin{lstlisting}[language=Python]
a = 10.0**16
print((a + 1.0) - a)
print(0.1 + 0.2 == 0.3)
\end{lstlisting}
\pause
\begin{block}{Typical binary64 output}
\texttt{0.0}\par\texttt{False}
\end{block}
The operations are rounded. The algebraic rewrite you expected need not preserve the computed answer.
\end{frame}
\section{Diagnose}
\begin{frame}{A model for one rounded operation}
For a normal result under standard assumptions,
\[
 \operatorname{fl}(a\circ b)=(a\circ b)(1+\delta),
 \qquad |\delta|\leq u.
\]
\small This relative-error model needs care near underflow and does not cover overflow, NaNs or division by zero. For binary64 round-to-nearest, $u=2^{-53}$.
\end{frame}
\begin{frame}{Separate the problem from the algorithm}
\begin{description}
\item[Conditioning] How much does the exact answer change under small input changes?
\item[Stability] Does the algorithm behave like a nearby well-posed computation?
\end{description}
\pause
\alert{A stable algorithm cannot remove intrinsic sensitivity in the problem.}
\end{frame}
\section{Repair}
\begin{frame}[label=exercise]{Find the dangerous subtraction}
For $ax^2+bx+c=0$,
\[
 x_{\pm}=\frac{-b\pm\sqrt{b^2-4ac}}{2a}.
\]
Take $a=c=1$ and $b=10^8$.
\begin{itemize}
\item Which branch subtracts nearly equal values?
\item How could the product $x_+x_-=c/a$ help?
\end{itemize}
\end{frame}
\begin{frame}{Use the large root to recover the small one}
For real roots, $a\ne0$, and nonzero $q$, set
\[
 q=-\tfrac12\left(b+\operatorname{sign}(b)\sqrt{b^2-4ac}\right),
 \quad x_1=q/a,\quad x_2=c/q.
\]
\small Take $\operatorname{sign}(0)=1$ in this formula. Handle repeated or zero roots separately. Scaling is still needed to avoid overflow in the discriminant.
\medskip
\hyperlink{exercise}{\beamerreturnbutton{Return to the exercise}}
\end{frame}
\begin{frame}[fragile]{Inspect a sum with an independent formulation}
\begin{lstlisting}[language=Python]
import math
xs = [1e16, 1.0, -1e16]
total = 0.0
for x in xs:
    total += x
print(total)         # 0.0
print(math.fsum(xs)) # 1.0
\end{lstlisting}
\small The exact real sum is 1. A changed summation method can preserve information lost by a left-to-right sum. Always state the arithmetic environment when comparing outputs.
\end{frame}
\section{Validate}
\begin{frame}{A residual is evidence about a stated equation}
Given a computed root $\widehat x$, inspect
\[
 r=a\widehat x^2+b\widehat x+c.
\]
\begin{itemize}
\item Scale the residual relative to the problem data.
\item Check a reference at higher precision.
\item Test limiting and adversarial inputs.
\end{itemize}
\end{frame}
\begin{frame}[fragile]{Leave a reproducible bug report}
\begin{alltt}
Inputs:      a = 1, b = 1e8, c = 1
Arithmetic:  binary64, nearest rounding
Observation: small root loses digits
Check:       product of roots and residual
\end{alltt}
\small Keep inputs, arithmetic assumptions and the independent check together. A screenshot of the final number is insufficient.
\end{frame}
\appendix
\begin{frame}[allowframebreaks]{Reading and lab checklist}
\small
\begin{thebibliography}{9}
\bibitem{goldberg} David Goldberg. \emph{What Every Computer Scientist Should Know About Floating-Point Arithmetic} (1991).
\href{https://docs.oracle.com/cd/E19957-01/806-3568/ncg_goldberg.html}{Author's article, hosted in Oracle's numerical guide}.
\end{thebibliography}
The examples and debugging sequence are original teaching material.
\framebreak
\begin{enumerate}
\item Reproduce the failure with the smallest input.
\item Distinguish representation, conditioning and algorithmic error.
\item Replace one numerically dangerous operation.
\item Validate using a different formulation.
\end{enumerate}
\end{frame}
\end{document}
